% Ad Atticum 1.1 --- headnote
% ad-atticum-01-01, July 65 BC, written at Rome

\textit{Cicero to Atticus, written at Rome in July 65 BC,
the year before his consular candidacy. The first surviving
letter on the consulship --- a long, candid analysis of
the field of competitors and the political ground. The
date is fixed by the reference to the tribunician elections
of 17 Sextilis (= 16 July; the canvassing for the
consulships of 64 BC was about to open). Publius Galba is
already in the field and being refused; Cicero takes that
as a sign of the strength of his own following. Of the
sure candidates, only Catiline is named with menace ---
``if it shall have been judged that the sun does not shine
at noon, he will be a certain competitor.'' (Catiline was
about to be tried for extortion in his governorship of
Africa; an acquittal, foreseen and bought, was the only
thing that would let him stand.) Cicero is plotting a
trip to Cisalpine Gaul as legate to Piso to court the
voting strength there, and asks Atticus to secure
Pompey's faction in his absence.}

\textit{The third paragraph is the letter's nervous
centre. Atticus's uncle Caecilius is suing Satyrus for
fraud, with the Luculli and Scipio joining the creditors.
Satyrus is one of Cicero's daily attendants, ranking only
behind Lucius Domitius (on whose canvassing-network
Cicero's whole consulship rests), and has rendered great
services to both Cicero and Quintus. Cicero has refused
to take Caecilius's brief, and Caecilius has taken the
refusal coldly. The whole paragraph is an apology to
Atticus for the refusal, ending on a quotation of
\textit{Iliad} 22.159 --- the runner Hector, racing
Achilles for his life, runs not for a sacrificial animal
or a hide but for a man's life --- with the political
analogue clear: this canvass is large.}

\textit{The closing line returns to the running theme of
the corpus: a Hermathena from Atticus has arrived and is
beautifully set up, ``so that the whole gymnasium seems
an offering to it.''}
